\documentclass[11pt,a4paper]{article}
\usepackage[utf8]{inputenc}
\usepackage[T1]{fontenc}
\usepackage[spanish]{babel}
\usepackage{hyperref}
\usepackage{geometry}
\geometry{margin=2.5cm}
\usepackage{parskip}

\title{074 ALAS- Apertura, Lengua, Acomodo, Sostenido}
\author{Alejandro Valenzuela Rivera\\\small Universidad Autónoma de Nuevo León}
\date{\today}
\begin{document}
\maketitle

\begin{quote}
\textbf{Sesgo declarado.} Este trabajo se sostiene desde el sesgo declarado de su autor, y no finge neutralidad. Lo que afirma se puede tumbar, y ese es el punto: el falsificador está escrito, no escondido.
\end{quote}

\section*{Resumen}
El estandar de protocolos de comunicacion: toda persona o entidad define su propio protocolo, y todo protocolo declara cuatro cosas. La apertura es desde donde se habla y en que modo; la lengua es en que se habla y que se traduce; el acomodo es converger o divergir y que se pierde; el sostenido es que puede el otro sin el que hablo.

\section{Por qué se sostiene}
Una comunicacion que funciona deja al otro con ALAS, no con el mensaje. La prueba no es que entienda: es que pueda seguir sin el que le hablo. Y las dos letras que arreglan el truncado son Lengua y Acomodo: un mensaje que no cabe en el canal no es claro, es truncado, y el truncado se lee como completo. Eso es lo que pasa con la gente que no entiende el idioma, que no oye, que no ve o que procesa distinto.

\section{De dónde viene}
Los tres protocolos de RUTA y la autoridad.

\section{El caso}
Alas: lo que se le deja al otro, no lo que se le dice.

\section{Método}
Un protocolo que declara apertura, lengua, acomodo y sostenido.

\section{Límites, y qué lo tumba}
Un protocolo sin lengua declarada, o con una acomodacion sin perdida declarada.

\section{Lo que queda abierto}
\textquestiondown{}Que declara un protocolo para no truncar a quien no es como el que habla?

\section*{Cómo se leyó esta tesis}
Esto se escribió en el repositorio man-made-empathy, del que este candado es una pieza: candados/06 del caso/074 ALAS- Apertura, Lengua, Acomodo, Sostenido.md. La máquina que lo compuso lo hizo desde este mismo archivo, y no agregó nada que no estuviera escrito. Si algo de esto es falso, se dice en el mismo sitio, y se corrige con fecha.

\end{document}
